\begin{tcolorbox}[colback=lightgray,colframe=deepnavy,arc=3pt,title=\textbf{RESUMO},fonttitle=\bfseries]
\small
Este trabalho apresenta o desenvolvimento parcial de uma aplicação para visualização de árvores arteriais geradas pelo método \emph{Constrained Constructive Optimization} (CCO), utilizando modelos geométricos fornecidos em formato VTK. O objetivo foi desenvolver uma aplicação de Computação Gráfica que integre leitura, representação, transformação, rasterização e recorte de modelos CCO, destacando quais algoritmos clássicos foram implementados pelo grupo e quais foram delegados à biblioteca. A aplicação foi desenvolvida em Python com PyVista e reúne leitura dos dados, representação dos segmentos vasculares por tubos, codificação visual do raio, controles de câmera e projeção, preenchimento de polígono pelo algoritmo Scanline e seleção de técnicas de antialiasing disponibilizadas pela biblioteca. O modelo tridimensional \texttt{tree3D\_Nterm016.vtk} foi utilizado como exemplo principal. Para esse arquivo, foram identificados 32 pontos, 31 células lineares e raios entre 0,2261241 e 0,6870107. O preenchimento Scanline aplicado a um polígono côncavo de teste ocupou 216 dos 780 pixels da grade (27,7\%), conforme indicado pela própria execução. Os módulos de rasterização de linhas e de recorte de segmentos e polígonos permanecem sem implementação; além disso, o módulo de câmera oferece vistas predefinidas, mas não implementa controles próprios de translação, rotação e escala. Assim, o relatório distingue os recursos efetivamente implementados daqueles previstos no enunciado e ainda pendentes.
\vspace{0.15cm}

\textbf{Palavras-chave:} Computação Gráfica; PyVista; CCO; rasterização; Scanline; visualização 3D.
\end{tcolorbox}